% 規劃修訂，尚未執行；不可作為既有第四章的實驗結果。
% 使用既有論文導言中的數學與 hyperref 套件。未修改控制器與既有章節。

\section*{方法導向實驗與消融規劃}

本版依「僅做模擬、著重方法」修訂原規劃。
模擬器是共同的測試環境，主要研究問題改為各項控制設計對追蹤、
避障與任務效率的影響。故障注入、模擬配速、物理引擎、
接觸模型切換及驅動增益掃描不列為主要比較。
必要的資料對齊與執行檢查仍在正式測試前完成，避免將實作錯誤帶入方法比較。
本文件是新實驗的工作規劃，不變更第四章已有數據。

\subsection*{一、核心比較與預算}

保留三層證據：外部方法比較、內部方法消融、完整任務驗證。
核心預算共 $792$ 次執行，包含 $780$ 趟正式任務測試與 $12$ 段模型資料蒐集。
模型資料蒐集又分為三段訓練、三段驗證與六段測試，
不能將全部 $792$ 次稱為獨立任務測試。
調參、網路訓練及開發試跑另計。
若加入選位、門扇位姿限制及完整時域曲線三項選配，另增加 $90$ 趟，合計 $882$ 次。

\begin{table}[htbp]
    \centering
    \caption{方法導向的核心實驗配置；消融只計新增執行}
    \label{tab:method_focused_budget}
    \small
    \renewcommand{\arraystretch}{1.2}
    \begin{tabular}{l p{0.64\textwidth} r}
        \hline
        編號 & 比較設計 & 次數 \\
        \hline
        M0 & 固定名目配置，三個構型區域，各四段獨立激勵紀錄 & $12$ \\
        N1 & 三類布局、靜態與交會兩條件、十種子、四導航方法 & $240$ \\
        N2 & 兩個交會場景、十種子、四個新增導航消融配置 & $80$ \\
        W1 & 四類操作、一般與受限空間、十種子、三策略 & $240$ \\
        W2 & 圓柱與抽屜、十種子、五個新增全身控制消融配置 & $100$ \\
        E0 & 兩住宅布局、無與有行人、十種子、三策略 & $120$ \\
        \hline
        合計 & $780$ 趟任務測試及 $12$ 段模型資料 & $792$ \\
        \hline
    \end{tabular}
\end{table}

導航保留完整方法、MPPI Omni、DWB Omni 及 RPP。
DWB 官方實作支援全向機器人；RPP 的側向運動能力差異須明列，
不將四組解讀為完全相同運動形式的純演算法比較。
參考：\url{https://api.nav2.org/nav2-jazzy/html/md_nav2_dwb_controller_README.html}。
四組採相同合成平移速度上限、輪速、命令變化限制與原始雷射防護，
並各自保存防護前後命令與介入時間。

W1 與 E0 統一使用三策略：本研究 WB-MPC、移植的 Holistic 反應式全身 QP，
以及本研究 WB-MPC 加停駐後接近。
Holistic 作者提供全向平台實作，仍需說明本研究對位置命令介面的移植方式，
不能直接宣稱重現原平台性能。
參考：\url{https://jhavl.github.io/holistic/}。
停駐後接近只要求底盤先到位停妥，再展臂、接近與建立抓持；
開門或拉抽屜期間仍可全身協同。
本版不另要求整段機構操作皆固定底盤的策略。

W1 四類操作為圓柱取放收納、筆狀物取放收納、門扇開關、抽屜開關。
所有比較從共同進場事件開始，計入方法自己的選位、接近、等待、操作與退出時間。
局部比較由共同 RGB-D 觀測建立初始世界座標目標及機構幾何，
同案例的各組使用同一目標封包；門與抽屜的後續目標沿相同機構模型更新。
E0 則保留線上視覺量測與重新觀測，評估整體閉迴路。

\subsection*{二、導航消融：保留五組，固定執行端防護}

\begin{table}[htbp]
    \centering
    \caption{導航方法消融配置；五組均保留相同的原始雷射防護}
    \label{tab:method_focused_nav_ablation}
    \small
    \renewcommand{\arraystretch}{1.2}
    \begin{tabular}{l c c c c}
        \hline
        配置 & 障礙運動預測 & CBF & 速度調整 & 側移調整 \\
        \hline
        A 完整方法 & 開 & 開 & 開 & 開 \\
        B 不使用障礙速度 & 關 & 開 & 開 & 開 \\
        C 移除 CBF & 開 & 關 & 開 & 開 \\
        D 移除參考調整 & 開 & 開 & 關 & 關 \\
        E 僅作速度調整 & 開 & 開 & 開 & 關 \\
        \hline
    \end{tabular}
\end{table}

A 與 B 比較障礙運動預測：保留同一追蹤器、物件關聯與目前幾何，
只在候選外推及 CBF 兩處將障礙速度設為零，不同時關閉雷射或追蹤器。
A 與 C 比較 CBF 的增量效果，移除對應限制列與無用的鬆弛成本，
其餘參考生成、輸入限制及防護保持一致。
D 與 E 比較速度候選，E 與 A 比較側移候選，D 與 A 則比較完整參考調整層。
這些成對比較可由五組結果取得，無須各自重跑完整方法。

使用開闊交會與狹窄交會兩種場景，每種十個不同案例。
五組共有 $100$ 趟資料，其中 A 的 $20$ 趟取自相同 N1 案例，
因此只新增 $80$ 趟。
固定障礙軌跡、起跑相位、速度上限與間距設定；
不以降低消融組的防護能力製造優勢。

主指標為完成率、實際非指定接觸、最小淨空、
共同成功案例的配對時間及等待時間。
另列雷射防護介入時長、介入比例與命令改動量，
以及 CBF 啟用時的鬆弛量。
若共同防護使兩組皆無接觸，仍可由介入量與時間觀察控制方法差異，
不能只以接觸皆為零判斷兩組相同。
本設計衡量保留共同防護時的模組效果，不宣稱關閉防護後仍有同樣結果。

\subsection*{三、全身控制消融：模型、預測與補償分開回答}

\begin{table}[htbp]
    \centering
    \caption{全身方法的核心消融配置}
    \label{tab:method_focused_wb_ablation}
    \small
    \renewcommand{\arraystretch}{1.2}
    \begin{tabular}{l c c c c}
        \hline
        配置 & 關節預測模型 & 預測步數 & 線上偏置 & TCP 積分 \\
        \hline
        A 完整方法 & 增廣 & $N_\star$ & 開 & 開 \\
        B 固定偏置 & 增廣 & $N_\star$ & 關 & 開 \\
        C 關閉 TCP 積分 & 增廣 & $N_\star$ & 開 & 關 \\
        D 兩種補償關閉 & 增廣 & $N_\star$ & 關 & 關 \\
        E 直接積分模型 & 直接積分 & $N_\star$ & 關 & 關 \\
        F 單步控制 & 增廣 & $1$ & 開 & 開 \\
        \hline
    \end{tabular}
\end{table}

D 與 E 回答設定點增廣模型是否有作用。
兩組皆關閉線上偏置與 TCP 積分，D 使用固定的離線辨識參數，
E 將設定點變化率直接積分為預測關節位置。
這是已辨識追蹤模型與直接積分假設的比較，不將差異全部歸因於單一參數。
不能在移除增廣模型的同時保留只有該模型才有定義的線上偏置機制，
再將結果稱為單因素比較。

A 與 F 回答多步預測相對單步控制的作用。
$N_\star$ 由開發集在 $\{5,10,20\}$ 中選定，
以相同控制週期、時間尺度一致的階段成本、固定終端矩陣、限制及參考生成器比較。
時域長度改變仍同時改變成本累加範圍與終端時刻，需一併說明。
$N=1$ 是本方法的內部消融，不能取代外部 Holistic 基線。

A、B、C、D 形成線上偏置與 TCP 積分的二乘二設計。
除 A 對 B、A 對 C，也檢查 C 對 D 與 B 對 D，
觀察兩種補償是否互補或遮蔽彼此效果。
固定偏置是保留共同離線辨識的初始值，並非任意設為零。
線上更新只使用當時已完成的追蹤區間，所有條件沿用相同更新門檻。
需報告偏置與積分各自的實際啟用時間；
若某案例未啟用，就不能用該案例宣稱補償有效。

使用圓柱取放與抽屜開關兩種任務，每種十個案例。
六組共 $120$ 趟，其中完整方法 A 的 $20$ 趟取自 W1 一般空間案例，
只新增 $100$ 趟。
採共同目標封包、機構模型、執行延遲處理、位置介面與接觸設定。
不另外掃描模擬驅動增益，也不刻意對某組加入較大的命令延遲。

指標包括 TCP 位置與姿態誤差、機構相容誤差、
關節與設定點限制餘量、實際放行命令變化、完成率、分段時間與求解耗時。
追蹤誤差使用原始操作目標，不只對補償後的目標計算，避免以修改目標降低表面誤差。
誤差依展開、接近、機構操作與運送分段，不能只報全程平均。

\subsection*{四、與貢獻主張連動的選配實驗}

若將底盤選位成本列為獨立貢獻，加入 S1：
四類操作各十案例，完整方法與最近可行底盤位置比較。
兩者共享候選生成、導航可達性、IK、關節與淨空檢查，只改變可行候選的排序。
完整方法的 $40$ 趟沿用 W1 一般空間資料，因此新增 $40$ 趟。
分析可行候選數、最小奇異值、關節餘量、完整路程、任務時間及成功率。

若將門扇硬位姿容許範圍列為獨立貢獻，加入 K1：
門扇十案例，僅移除硬位姿限制列，保留同一機構參考、權重與接觸設定。
新增 $10$ 趟，對照 W1 的十趟完整方法，分析最大相容誤差與可行性代價。

若需完整時域敏感度曲線，加入 H1：
在 W2 的二十個案例補齊 $\{5,10,20\}$ 中另外兩個步數，新增 $40$ 趟。
與既有 $N_\star$ 及 $N=1$ 一起呈現精度、命令變化與計算量的關係。
這些是加做的配置，不能只挑測試結果最好的步數作為主方法。

\subsection*{五、不列為主要實驗的內容}

\begin{itemize}
    \item 不比較 Isaac Sim 與 Gazebo 哪個較好；各項方法在同一測試平台中比較。
    \item 不以物理步長、模擬配速、CPU/GPU 型號或通訊排程作為獨立實驗因素。
    \item 不安排主動注入求解逾時、程序故障及封包失效的正式實驗組。
    \item 不做固定連接與摩擦接觸的模型對照，也不掃描材質、摩擦或驅動增益。
    \item 不將現成 Mask R-CNN 的網路架構消融列為控制方法的核心實驗。
    感知只需交代固定模型、獨立測試誤差與使用限制；E0 保留實際線上感知。
    \item 在途命令前推與停止備援仍維持一致，作為共同執行條件，
    本版不將其獨立效益列為主要實驗主張。
\end{itemize}

不另做接觸模型比較，仍須如實記錄使用的接觸設定。
若物件以固定相對連接跟隨末端，結果可支持運動學協同及流程銜接，
不能據此推論摩擦抓持、滑移抵抗或真實接觸可靠性。
本版不要求為了比較控制方法另外開發新的接觸模型。

\subsection*{六、完整住宅任務與結果判讀}

E0 保留 $120$ 趟：兩個布局、無行人與有行人、十個種子及三策略。
第二布局不參與調參；行人軌跡及互動規則事前固定。
完整流程保留開門、通行、圓柱與筆收納、關抽屜、回程關門及返回。
WB-MPC 與 Holistic 的比較是完整控制配置比較；
WB-MPC 與停駐後接近的比較則檢查移動與操作重疊的效益。
不能將這兩種對照混為相同的單因素消融。

不同種子需改變保存的初始構型、物件位置或交會相位，
不能將同一確定性案例重播十次視為十個獨立場景。
同案例的方法與消融配置放在同一執行區塊，使用預先保存的隨機順序。
共用完整方法資料須同時滿足程式版本、資產、目標封包、參數、
完成條件與紀錄方式一致；不一致時須重跑，並重新計算新增預算。

十個種子是各分層的起始工作預算，不保證能辨識很小的差異。
先以開發案例估計配對差的變異與失敗情形，正式測試前固定樣本數及停止規則，
若需增加則對所有比較組共同增加，不只補跑不理想的結果。
成功率呈現區間，時間同時交代兩組皆成功的配對數，
並對所有嘗試呈現失敗與逾時，避免只比較容易成功的子集合。

各案例先取得分段統計，再以案例為配對單位呈現差值與信賴區間，
不將每個控制步當作獨立樣本。
預先指定主要對照與指標，多組比較時控制多重比較問題。
若某消融沒有明顯影響，需據實說明其作用條件，
或將該項由獨立貢獻降為實作選項，不預先承諾完整方法必定全面勝出。

\subsection*{七、預計形成的主要圖表}

\begin{enumerate}
    \item 導航外部基線：各布局完成率、淨空分布、配對時間與防護介入量。
    \item 導航五組消融：各主要對照的效果量與信賴區間，配合交會軌跡及速度曲線。
    \item 全身外部基線：四類操作的完成率、TCP 分段誤差與底盤手臂運動分工。
    \item 全身六組消融：模型對照、單步與多步對照，以及補償二乘二效果圖。
    \item 住宅整體任務：各子任務失敗分布、總時間與階段時間，並呈現失敗案例。
    \item 若加做 S1、K1 或 H1，再分別提供選位、機構誤差或時域敏感度圖。
\end{enumerate}

本文件與附帶 CSV 僅建立規劃，沒有啟動模擬、訓練或控制器執行。
